\documentclass[11pt,a4paper]{article}
\usepackage[utf8]{inputenc}
\usepackage[T2A]{fontenc}
\usepackage[russian]{babel}
\usepackage{amsmath,amssymb,amsfonts,mathtools}
\usepackage{booktabs}
\usepackage{geometry}
\usepackage{hyperref}
\usepackage{array}
\geometry{left=2.5cm,right=2.5cm,top=2.5cm,bottom=2.5cm}

\title{Единая геометродинамическая модель элементарных частиц\\
на основе дислокационных колец:\\
кулон, гравитация и мёбиусова частица}
\author{А.~Дроздов}
\date{23 октября 2026}

\begin{document}
\maketitle

\begin{abstract}
Излагается единая геометродинамическая модель, в которой электрический
заряд, гравитация и спин элементарных частиц возникают из геометрии
дислокационных колец в калибровочной теории дислокаций Лазара--Брагинского~\cite{LazarDyn2010,Braginsky2018}
и динамики скольжения трёхмерной упругой <<мыльной плёнки>> --- границы
четырёхмерного интеграла действия. Электрон моделируется как два
прецессирующих дислокационных кольца (диполь), поле строится из двух
каналов: центрального (кулоно-аналог) и вихревого (аналог магнитного поля,
сила Пича--Кёлера). Показано, что гравитационно-подобное притяжение
получается в правильной функциональной форме $1/R^2$ с верным знаком, а
введённая автором \emph{мгновенная нелокальная плёнка} устраняет разрушение
закона в дальней зоне (численно $F\propto\cos\Delta\varphi/R^2$, $K\to4$).
Установлено тождество $(\ell_P/a_{comp})^2=\alpha_G$ и доказано четырьмя
независимыми путями, что величина $G$ не выводится из постоянных
$(\hbar,c,e)$: гравитации необходим отдельный планковский сектор.
Отдельно исследована мёбиусова топология контура как модель частицы,
отличной от электрона (период $4\pi$, дальнодействие $1/R$ от
топологической особой точки, класс майорана-подобного фермиона).
\end{abstract}

\tableofcontents

\section{Введение}
Современная физика разделена барьером между Общей теорией относительности
и квантовой механикой. Традиционные попытки квантования гравитации
сталкиваются с ультрафиолетовыми расходимостями, вызванными абстракцией
точечных частиц. В настоящей работе предлагается альтернативный подход,
восходящий к идее Э.~Ферми (1923) о геометрии интегрирования упругих полей~\cite{Fermi1923}
и к калибровочной теории дислокаций в кристаллах, где заряды возникают как
топологические дефекты решётки~\cite{KadicEdelen1983,LazarDyn2010}. Все взаимодействия трактуются геометрически:
как динамика упругой границы интеграла действия вдоль мировых трубок
прецессирующих дефектов, а не как силовые поля или искривление
пространства-времени.

\section{Поля Лиенара--Вихерта и поправка Ферми}
Базовое поле совпадает с формулой (63,8) Ландау--Лифшица~\cite{LL2}. Суммарное
электрическое поле движущегося заряда:
\begin{equation}
\vec{E}=\frac{q}{R^{*3}}\left[\left(\vec R-\frac{R}{c}\vec v\right)
\left(1+\frac{\vec R\vec a}{c^2}-\frac{v^2}{c^2}\right)
-\vec a\,\frac{R^{*}R}{c^2}\right],\qquad R^{*}=R-\frac{\vec v\vec R}{c}.
\end{equation}
Тождественность формуле (63,8) проверяется раскрытием двойного векторного
произведения ускорительного члена по правилу $bac-cab$ с использованием
$\vec R\cdot(\vec R-\tfrac{R}{c}\vec v)=R\,R^{*}$. Авторское дополнение ---
множитель Ферми $\left(1+\vec R\vec a/c^2\right)$, отвечающий искривлению
мировой трубки при ускорении (вариация действия <<типа B>>).

\section{Модель электрона: два дислокационных кольца}
Электрон представляется двумя прецессирующими дислокационными кольцами
(диполь $+q_m,-q_m$). Поле элемента кольца раскладывается на два канала:
\begin{equation}
\text{central}\ (\text{кулон}),\qquad \text{vortex}=\vec n\times\text{central}\ (\text{Пич--Кёлер}).
\end{equation}
Поле всегда запаздывающее (досветовое); поправка Ферми реализована в трёх
режимах: PLAIN, FILM\_RETARDED и FILM\_INSTANT.

\section{Мгновенная нелокальная плёнка}
Ключевой элемент модели требует точного разграничения двух уровней.
\textbf{(1) Поле.} Электромагнитное (дислокационное) поле --- результат
минимизации \emph{самого} интеграла действия: уравнения типа Максвелла,
потенциалы Лиенара--Вихерта, силы Кулона и Пича--Кёлера~\cite{LazarDyn2010,Braginsky2018}.
Это поле \emph{запаздывает по классике и содержит радиационный член};
излучение не подавляется. \textbf{(2) Поправка Ферми.} Трёхмерная упругая
граница 4-интеграла действия (<<мыльная плёнка>>) подстраивается
\emph{мгновенно и нелокально} (эллиптическая задача, не волновая),
скользя по мировым трубкам с сохранением ортогональности (условие Ферми) и
минимизируя поверхностную энергию. Результат этой отдельной минимизации ---
поправка к дифференциалу интервала $ds$ --- и есть то единственное, что
распространяется без запаздывания; у неё нет собственного масштаба длины
волны, поэтому она не разрушается в дальней зоне. Таким образом
мгновенность относится к наклону плёнки, а не к полю: плёнка не запрещает
излучение, а лишь мгновенно задаёт границу интегрирования.

Для численной оценки эффекта поправки использован безызлучательный
модельный прокси (аналог Веберовской электродинамики: скоростной вклад без
радиационного члена) --- не как утверждение о структуре поля, а как
вычислительное упрощение для изоляции вклада мгновенного наклона. В этом
пределе установлено:
\begin{equation}
F(R,\Delta\varphi)\approx-K\left(\frac{v}{c}\right)^2
\frac{\cos\Delta\varphi}{R^2},\qquad K\to4\ (v/c\to0),
\end{equation}
устойчиво на диапазоне $R/a=3\ldots8000$ без смен знака, тогда как
запаздывающий механизм разрушает закон при $R\gtrsim c/\omega$.

\section{Три варианта кулона и гравитации}
\begin{center}
\begin{tabular}{clcc}
\toprule
Вариант & Механизм & Порядок & Итог \\
\midrule
1 & линейная магнитостатика & $R^{-3}$ & ни кулон, ни гравитация \\
2 & изотропный источник Ферми $\sigma=mW^2/c^2\ge0$ & $R^{-1}$ & гравитация (форма $\checkmark$, величина $\times$) \\
3 & нелинейная самосвязь (Янг--Миллс) & $R^{-1}\cos\Delta\varphi$ & кулон (форма $\checkmark$, модуль $f_1$ свободен) \\
\bottomrule
\end{tabular}
\end{center}
Гравитационный канал (Вариант~2) даёт истинно притягивающее поле
($\sigma\ge0$ независимо от знака и фазы), размерностно чистое:
$\delta w=-\sigma/(\gamma_{3D}R)$ безразмерно, $\gamma_{3D}=c^7/(8\pi\hbar G^2)$.

\section{Барьер величины $G$: необходимость планковского сектора}
Установлено тождество
\begin{equation}
\left(\frac{\ell_P}{a_{comp}}\right)^2=\frac{Gm_e^2}{\hbar c}=\alpha_G
\approx1{,}75\times10^{-45},
\end{equation}
проверенное с точностью $1{,}8\times10^{-16}$. Оно переводит вопрос о
малости гравитации в вопрос <<почему $b=\ell_P$>>. Четыре независимых пути
вывести $b=\ell_P$ без явного ввода $G$ дают отрицательный результат:
\begin{enumerate}
\item тождество масштаба --- круговая зависимость ($G$ в $\ell_P$);
\item энергетическое замыкание $m_ec^2=\gamma_{3D}b^2a$ --- мимо на $\sim10^{44}$;
\item условие устойчивости --- $(b/\ell_P)^2=4\alpha\,\alpha_G$;
\item квантование потока (Брагинский~\cite{Braginsky2018}) --- квант на масштабе $(\hbar,c,e)$, не $\ell_P$.
\end{enumerate}
\textbf{Структурный вывод.} Квант Бюргерса, фиксируемый постоянными
$(\hbar,c,e)$, несовместим с планковской длиной. Кулон (модуль $a_1$) и
гравитация (модуль $f_1$) управляются разными квантами; $G$ требует
отдельного, независимо постулируемого планковского сектора. Одновременное
описание кулона и гравитации остаётся непротиворечивой двухмодульной
калибровкой, но не предсказанием величины $G$.

\section{Устойчивость дислокационного кольца}
Устойчивость обеспечивается балансом кулоновской самоэнергии и линейного
натяжения дислокации $T=\mu b^2$:
\begin{equation}
E(a)=\frac{kq^2}{\pi a}\ln\frac{8a}{r_0}+T\cdot2\pi a,\qquad
a_{eq}\simeq\sqrt{\frac{kq^2}{2\pi T}},\quad E''(a_{eq})>0.
\end{equation}
Численно подтверждён устойчивый минимум и скейлинг $a_{eq}\propto1/\sqrt{T}$.

\section{Мёбиусова частица (не электрон)}
Мёбиусова топология контура допускает два уровня: Уровень~2 ---
переформулировка знаковой пары электрона, Уровень~1 --- самостоятельная
частица. Для одиночного мёбиусова кольца знак заряда несёт
антипериодическая структура $\cos(\theta/2)$ (период $4\pi$, спинорная
двузначность). Численно: монопольный момент равен нулю, но поле убывает как
$1/R$ ($p=1{,}00$ устойчиво на четырёх порядках по $R$) --- следствие
локализованной топологической особенности (<<шва>>), а не электрического
монополя. Особая точка регуляризуется как топологический кинк с конечной
энергией $E_{core}=2\sqrt{\kappa T}$, что даёт конечную массу
$m=2\mu b^2 a/c^2$ и топологическую стабильность. По набору квантовых чисел
(нейтральность, геометрический спин-$\tfrac12$, самосопряжённость) ближайший
аналог --- майорана-подобный фермион, а не частица Стандартной модели.

\section{Парадокс 4/3 для настоящих частиц и вариация B Ферми}
\label{sec:four-thirds}

Классический парадокс 4/3 (расхождение между электромагнитной массой из
энергии $m_U=U/c^2$ и из импульса поля $m_p=p/v$) допускает перенос на
упругие дефекты. Множитель $k=m_p/m_U$ зависит от типа дефекта.

\paragraph{Винтовая компонента.} Для винтовой дислокации поле антиплоское
(эффективно скалярное, одна поперечная мода $c_t$). Движущаяся винтовая
дислокация ведёт себя как релятивистская струна,
$E(v)=E_0/\sqrt{1-v^2/c_t^2}$~\cite{Eshelby1949,HirthLothe}, откуда
$k_{\text{screw}}=1$ --- парадокса 4/3 нет.

\paragraph{Краевая компонента.} Краевая дислокация вовлекает обе моды
($c_t$ и $c_l$), $(c_t/c_l)^2=(1-2\nu)/(2(1-\nu))$. В ведущем
низкоскоростном порядке эффективный множитель
\begin{equation}
k_{\text{edge}}(\nu)\simeq\frac{1}{1-\nu},
\end{equation}
растущий от $1$ при $\nu=0$ до $2$ при $\nu\to1/2$. Примечательно, что при
условии Коши $\nu=1/4$ (центрально-силовой изотропный континуум ---
каноническая модель <<эфира>>, $c_l=\sqrt{3}\,c_t$) получается
\begin{equation}
k_{\text{edge}}(1/4)=\frac{1}{1-1/4}=\frac{4}{3},
\end{equation}
то есть ровно электродинамическое значение. Это даёт гипотезу: множитель
$4/3$ Максвелла может быть проявлением упругого $4/3$ коши-эфира.

\paragraph{Снятие 4/3 вариацией B Ферми (тождество с рецептом Белинфанте).}
Парадокс 4/3 --- артефакт интегрирования $\int T^{0i}$ по гиперплоскости
лабораторной одновременности~\cite{Rohrlich}. При интегрировании по
гиперплоскости $\Sigma$, ортогональной 4-скорости $u^\mu$ дефекта,
$P^\mu=\int_\Sigma T^{\mu\nu}d\Sigma_\nu=(U/c^2)u^\mu$ и $k=1$. В формализме
Лазара~\cite{LazarDyn2010} 4-импульс дефекта $P_a=\int_\Sigma(\pi_a+p_a)$ с
дислокационным импульсом $\pi_a=(e_a\lrcorner T^b)\wedge D_b$; при $\Sigma$,
нормальной мировой трубке, поток дислокационного напряжения $E_a$ через
границу обнуляется (полевое уравнение $dH_a-\dot D_a-E_a=\Sigma_a$), что
даёт $k=1$. Условие <<$\Sigma$ нормальна мировой трубке>> есть дословно
вариация B Ферми. Таким образом \textbf{вариация B Ферми тождественна
ковариантному рецепту Белинфанте}, и обе снимают 4/3 --- как для электрона,
так и для настоящих частиц (атомов эфира с краевой компонентой).

\section{Заключение}
Модель даёт связную геометрическую картину, объясняющую функциональную
\emph{форму} кулона, гравитации и спина из единого источника, и предлагает
проверяемый элемент --- мгновенную нелокальную плёнку, снимающую проблему
дальней зоны. Решены: форма $1/R^2$, знак, размерности, устойчивость кольца
и стабильность мёбиусовой частицы. Единственный принципиальный барьер ---
величина $G$ --- локализован точно: она требует отдельного планковского
сектора, не выводимого из $(\hbar,c,e)$, что доказано четырьмя независимыми
путями. Масса мёбиусовой частицы и количественная связь кулон--гравитация
упираются в тот же корень.

\paragraph{Открытые вопросы.}
(i) обоснование планковского сектора без явного ввода $G$;
(ii) совместимость неориентируемого мёбиусова контура с формализмом Лазара;
(iii) строгий вывод мгновенной плёнки как PDE минимальной гиперповерхности
для $N$ тел;
(iv) вывод абсолютной массы мёбиусовой частицы из натяжения и ширины кинка.

\paragraph{Замечание о воспроизводимости.}
Все численные результаты получены независимым Python-портом эталонного
C-кода \texttt{dislocation\_rings\_v4}; формулы и параметры взяты из
исходных файлов проекта.



\appendix

\section{Механизм Хиггса и ядро дислокации (по Lazar 2009)}
\label{app:higgs}
Работа Lazar~\cite{LazarHiggs2009} даёт теоретическую опору под конечность
ядра мёбиусовой особой точки (раздел~\ref{app:num}, кинк). В калибровочной
теории дислокация --- не $\delta$-струна на линии с искусственной
расходимостью, а <<размазанный>> топологический объект; ядро возникает
естественно, и искусственный радиус обрезания становится излишним. При
минимальной связи дислокационной теории с асимметричной упругостью
калибровочно-инвариантная дисторсия $\beta_{ij}=\nabla_j u_i=u_{i,j}+\phi_{ij}$
поглощает намбу-голдстоуновское поле $u_i$ и приобретает <<массу>> ---
это и есть механизм Хиггса: $\beta_{ij}$ становится полем Прока с конечной
обратной длиной (массой), задающей масштаб ядра. Отсюда вытекает, что
ширина кинка $w$ в нашей модели (раздел~\ref{app:num}) не свободный
параметр, а внутренняя длина теории, определяемая материальными
постоянными $c_1,c_2,c_3,\mu,\gamma$; Lazar и Анастассиадис находят
несколько таких характерных длин. Это намечает путь к выводу абсолютной
массы мёбиусовой частицы --- через прока-массу поля $\beta_{ij}$, а не
через произвольную обрезку.

\section{Численные результаты}
\label{app:num}
Все значения получены независимым Python-портом эталонного C-кода
\texttt{dislocation\_rings\_v4}; $c=1$, $a$ --- радиус кольца.

\paragraph{Б.1. Мгновенный наклон плёнки (безызлучательный прокси).}
Коэффициент $K$ в $F\approx-K(v/c)^2\cos\Delta\varphi/R^2$ по мере роста
$R/a$ (сходимость к $K\to4$):
\begin{center}
\begin{tabular}{rccc}
\toprule
$R/a$ & $v/c=0.1$ & $v/c=0.03$ & $v/c=0.01$ \\
\midrule
3    & 724.6 & 8166 & 73589 \\
100  & -3.76 & -1.33 & 20.0 \\
1000 & -3.998 & -3.973 & -3.76 \\
8000 & -4.000 & -4.000 & -3.996 \\
\bottomrule
\end{tabular}
\end{center}

\paragraph{Б.2. Мёбиусова частица: асимптотика поля.}
Потенциал $V(R)$ одиночного мёбиусова кольца; плато $V\!\cdot\!R$
подтверждает закон $1/R$ ($p=1.00$):
\begin{center}
\begin{tabular}{rcc}
\toprule
$R/a$ & $V\cdot R$ & показатель $p$ \\
\midrule
100   & $1.587\times10^{-4}$ & --- \\
1000  & $1.572\times10^{-4}$ & --- \\
30000 & $1.571\times10^{-4}$ & $1.002$ \\
\bottomrule
\end{tabular}
\end{center}
Регуляризация разрыва кинком $\tanh((\theta-\pi)/\delta)$ даёт $p=1.02$
устойчиво при $\delta=0.5,0.2,0.05$ --- хвост $1/R$ физичен.

\paragraph{Б.3. Устойчивость кольца (натяжение).}
Равновесный радиус $a_{eq}$ от натяжения $T$; $a_{eq}\sqrt{T}\approx$const
подтверждает $a_{eq}\propto1/\sqrt{T}$:
\begin{center}
\begin{tabular}{rcc}
\toprule
$T$ & $a_{eq}$ & $a_{eq}\sqrt{T}$ \\
\midrule
0.02 & 4.927 & 0.697 \\
0.05 & 3.036 & 0.679 \\
0.20 & 1.456 & 0.651 \\
0.50 & 0.893 & 0.632 \\
\bottomrule
\end{tabular}
\end{center}

\paragraph{Б.4. Барьер величины $G$: четыре пути.}
\begin{center}
\begin{tabular}{lcc}
\toprule
Путь & Результат для $b$ & $b=\ell_P$? \\
\midrule
Тождество масштаба       & $(\ell_P/a)^2=\alpha_G$ & круговая зависимость \\
Энергетическое замыкание & $b\sim10^{-57}$~м & нет ($\sim10^{44}$) \\
Устойчивость             & $(b/\ell_P)^2=4\alpha\alpha_G$ & нет \\
Квантование потока       & масштаб $(\hbar,c,e)$ & нет ($10^{22}$) \\
\bottomrule
\end{tabular}
\end{center}

\begin{thebibliography}{9}
\bibitem{Fermi1923} E.~Fermi, \emph{Correzione di una contraddizione tra la
teoria elettrodinamica e quella relativistica delle masse elettromagnetiche},
Il Nuovo Cimento \textbf{25}, 159--170 (1923); перепечатано в
\emph{Enrico Fermi: Note e Memorie (Collected Papers)}, т.~I,
Univ. of Chicago Press (1962), с.~26--32.
\bibitem{LL2} Л.~Д.~Ландау, Е.~М.~Лифшиц, \emph{Теория поля}
(Курс теоретической физики, т.~II), Наука, Москва --- формула~(63,8).
\bibitem{LazarDyn2010} M.~Lazar, \emph{On the fundamentals of the
three-dimensional translation gauge theory of dislocations},
arXiv:1003.3549 [cond-mat.mtrl-sci] (18 Mar 2010); динамическая версия
$T(3)$-калибровочной теории дислокаций.
\bibitem{LazarStatic2008} M.~Lazar, C.~Anastassiadis, \emph{The
$T(3)$-gauge theory of static dislocations: force and pseudomoment
stresses}, arXiv:0802.0670 [cond-mat.mtrl-sci] (2008).
\bibitem{LazarHiggs2009} M.~Lazar, \emph{On the Higgs mechanism and stress
functions in the translational gauge theory of dislocations},
arXiv:0903.0990 [cond-mat.mtrl-sci] (2009); ядро дислокации возникает
естественно, без искусственного обрезания.
\bibitem{KadicEdelen1983} A.~Kadi\'c, D.~G.~B.~Edelen, \emph{A Gauge Theory
of Dislocations and Disclinations}, Lecture Notes in Physics \textbf{174},
Springer (1983).
\bibitem{Braginsky2018} А.~Я.~Брагинский, \emph{Квантование вектора Бюргерса
через минимальное взаимодействие и динамическая сила Пича--Кёлера}
(2016/2018).
\bibitem{Eshelby1949} J.~D.~Eshelby, \emph{Uniformly moving
dislocations}, Proc.~Phys.~Soc.~A \textbf{62}, 307 (1949).
\bibitem{HirthLothe} J.~P.~Hirth, J.~Lothe, \emph{Theory of Dislocations},
2nd ed., Wiley (1982), гл.~7 (динамика дислокаций, движущаяся краевая).
\bibitem{Rohrlich} F.~Rohrlich, \emph{Classical Charged Particles},
Addison-Wesley (1965); о 4/3 как артефакте нековариантного интегрирования.
\end{thebibliography}

\end{document}
